% Generated from paper/full_draft. Edit the Markdown source or migration script.

\section{Conservative Baseline and Seed Stability Diagnostics}
\label{sec:12-appendix:conservative-baseline-and-seed-stability-diagnostics}

The conservative confidence-only context policy is the stable baseline for context compaction. The main text reports calibrated token-budget policies as the primary cost result; this appendix keeps the earlier confidence-only scale table and seed diagnostics. These intervals are engineering stability diagnostics, not strong inferential proof: each scale has only three observations.

\begin{table*}[tbp]
\centering
\small
\caption{Multi-seed retrieval-quality stability}
\label{tab:a1}
\begin{tabularx}{\textwidth}{>{\raggedright\arraybackslash}X>{\raggedleft\arraybackslash}X>{\raggedleft\arraybackslash}X>{\raggedleft\arraybackslash}X>{\raggedleft\arraybackslash}X>{\raggedleft\arraybackslash}X}
\toprule
Scale & Dense $\mathrm{Hit@10}$ & Policy $\mathrm{Hit@10}$ mean & Std & 95\% CI & Mean hit delta \\
\midrule
100k & 0.8674 & 0.8652 & 0.0035 & [0.8565, 0.8739] & -0.0022 \\
200k & 0.7970 & 0.8249 & 0.0079 & [0.8052, 0.8446] & +0.0280 \\
400k & 0.7718 & 0.7819 & 0.0044 & [0.7709, 0.7929] & +0.0101 \\
638k & 0.7282 & 0.7466 & 0.0089 & [0.7246, 0.7687] & +0.0185 \\
\bottomrule
\end{tabularx}
\end{table*}

\begin{table*}[tbp]
\centering
\small
\caption{Multi-seed final context-token stability}
\label{tab:a2}
\begin{tabularx}{\textwidth}{>{\raggedright\arraybackslash}X>{\raggedleft\arraybackslash}X>{\raggedleft\arraybackslash}X>{\raggedleft\arraybackslash}X>{\raggedleft\arraybackslash}X>{\raggedleft\arraybackslash}X>{\raggedleft\arraybackslash}X}
\toprule
Scale & Dense $\mathrm{Tokens@10}$ & Policy $\mathrm{Tokens@10}$ mean & Std & 95\% CI & Mean token saving & Saving 95\% CI \\
\midrule
100k & 1472.39 & 1401.24 & 11.49 & [1372.70, 1429.79] & 4.83\% & [2.89\%, 6.77\%] \\
200k & 1444.12 & 1376.46 & 4.61 & [1365.01, 1387.91] & 4.69\% & [3.89\%, 5.48\%] \\
400k & 1482.30 & 1403.43 & 31.10 & [1326.16, 1480.69] & 5.32\% & [0.11\%, 10.53\%] \\
638k & 1525.62 & 1451.49 & 3.83 & [1441.97, 1461.00] & 4.86\% & [4.24\%, 5.48\%] \\
\bottomrule
\end{tabularx}
\end{table*}

The token-saving direction is consistent across scales. The wider 400k token interval should be interpreted as route-confidence and context-budget variance, not hidden as a uniform result.

A five-seed extension is also available for LoTTE 100k. The additional seeds test sensitivity to KMeans initialization and route-policy stochasticity.

\begin{table*}[tbp]
\centering
\small
\caption{Five-seed LoTTE 100k extension}
\label{tab:a3}
\begin{tabularx}{\textwidth}{>{\raggedright\arraybackslash}X>{\raggedleft\arraybackslash}X>{\raggedleft\arraybackslash}X>{\raggedleft\arraybackslash}X>{\raggedleft\arraybackslash}X>{\raggedleft\arraybackslash}X}
\toprule
Setting & Seeds & $\mathrm{Hit@10}$ & Avg tokens@10 & Token ratio & Token saving \\
\midrule
Dense-only & 1 & 0.8674 & 1472.39 & 1.0000x & 0.00\% \\
Conservative policy & 5 & 0.8708 & 1399.83 & 0.9507x & 4.93\% \\
\bottomrule
\end{tabularx}
\end{table*}

The five-seed hit-delta interval is $[-0.0082, +0.0150]$. This supports dense-level retrieval quality with stable context-token reduction, not a statistical-superiority claim.

\section{Static Retrieval Baselines}
\label{sec:12-appendix:static-retrieval-baselines}

Dense retrieval is the primary quality baseline. BM25 supplies lexical coverage, but it is weaker as a standalone retriever on LoTTE. Static dense-plus-BM25 reciprocal-rank fusion is competitive at some scales but does not consistently dominate dense.

\begin{table*}[tbp]
\centering
\small
\caption{Static LoTTE retrieval baselines across corpus scale}
\label{tab:b1}
\begin{tabularx}{\textwidth}{>{\raggedright\arraybackslash}X>{\raggedleft\arraybackslash}X>{\raggedleft\arraybackslash}X>{\raggedleft\arraybackslash}X>{\raggedleft\arraybackslash}X}
\toprule
Scale & Corpus chunks & BM25 $\mathrm{Hit@10}$ & Dense $\mathrm{Hit@10}$ & Static hybrid $\mathrm{Hit@10}$ \\
\midrule
100k & 101311 & 0.7232 & 0.8674 & 0.8624 \\
200k & 201010 & 0.6292 & 0.7970 & 0.8003 \\
400k & 400674 & 0.5721 & 0.7718 & 0.7617 \\
638k & 638509 & 0.5084 & 0.7282 & 0.7181 \\
\bottomrule
\end{tabularx}
\end{table*}

The declining dense score as corpus scale grows motivates adaptive context control, but it does not make dense retrieval obsolete. Dense remains an important recall floor and fallback route in IntentWeight.

\section{Cost Metric Separation}
\label{sec:12-appendix:cost-metric-separation}

The experiments separate three efficiency layers:

\begin{enumerate}
\item source candidate cost: candidates considered before final fusion;
\item dense invocation rate: fraction of queries using global dense retrieval;
\item final context tokens: retrieved chunk tokens sent to the generator.
\end{enumerate}

The main paper claim uses the third layer. Historical routing experiments showed that reducing candidate counts does not automatically reduce final context tokens when the final context remains fixed at top-10.

\begin{table*}[tbp]
\centering
\small
\caption{Representative fixed-top-10 correction audit}
\label{tab:c1}
\begin{tabularx}{\textwidth}{>{\raggedright\arraybackslash}X>{\raggedright\arraybackslash}X>{\raggedleft\arraybackslash}X>{\raggedleft\arraybackslash}X>{\raggedleft\arraybackslash}X>{\raggedleft\arraybackslash}X}
\toprule
Dataset / scale & Routing setting & $\mathrm{Hit@10}$ & Avg tokens@10 & Ratio vs dense & Source candidate cost \\
\midrule
Banking77 & Gated cost-aware routing & 0.9813 & 120.82 & 0.9978x & 142.51 \\
eManual & Gated cost-aware routing & 0.0116 & 17.92 & 0.9829x & 214.07 \\
LoTTE 100k & Quality-first routing & 0.8770 & 1518.44 & 1.0313x & 229.97 \\
LoTTE 100k & Conditional fallback routing & 0.8747 & 1516.24 & 1.0298x & 227.29 \\
LoTTE 100k & Cluster-credit routing & 0.8764 & 1550.65 & 1.0532x & 181.47 \\
LoTTE 200k & Initial gated routing & 0.8154 & 1549.39 & 1.0729x & 232.01 \\
LoTTE 400k & Initial gated routing & 0.7836 & 1547.66 & 1.0441x & 233.22 \\
LoTTE 638k & Initial gated routing & 0.7343 & 1599.95 & 1.0487x & 236.22 \\
\bottomrule
\end{tabularx}
\end{table*}

This audit motivated the explicit confidence-based final context policy. Its token savings come from reducing selected final context size in high-confidence cases, not from treating candidate-count savings as prompt token savings.

\section{Secondary Datasets and Boundary Cases}
\label{sec:12-appendix:secondary-datasets-and-boundary-cases}

The secondary datasets have different roles and should not be pooled into the main LoTTE evidence claim.

\begin{itemize}
\item \textbf{PubMedQA} is an evidence-retrieval proof-of-concept near a dense ceiling. Dense reaches $\mathrm{Hit@10}=0.9930$; trust-weighted feedback reaches $\mathrm{Hit@10}=0.9940$, last reward $0.8727$, and selected-cluster hit $0.8860$. The ground truth is abstract-level context, not a strict answer sentence.
\item \textbf{Banking77} is an intent-routing proxy rather than an evidence-retrieval benchmark. Dense/reference $\mathrm{Hit@10}$ is $0.9805$; trust-weighted feedback reaches $\mathrm{Hit@10}=0.9844$, last reward $0.9805$, and selected-cluster hit $0.9983$. It supports the feedback mechanism, not the main evidence-retrieval headline.
\item \textbf{eManual} is a duplicate-text limitation case. Strict dense $\mathrm{Hit@10}$ is $0.3231$, text-equivalent dense $\mathrm{Hit@10}$ is $0.5615$, and deduplicated dense $\mathrm{Hit@10}$ rises to $0.8615$. Strict chunk IDs can therefore understate useful retrieval.
\item \textbf{CUAD} is a sparse legal smoke/stress case. Dense/reference $\mathrm{Hit@10}$ is $0.0759$ and the trust-weighted smoke reaches $\mathrm{Hit@10}=0.0886$. It is a GT-anchored sample, not positive full-corpus evidence.
\end{itemize}

\subsection{eManual Duplicate-Text Diagnostic}
\label{sec:12-appendix:emanual-duplicate-text-diagnostic}

eManual contains 18,812 corpus chunks but only 1,729 unique text strings. Strict chunk-id evaluation can therefore mark semantically equivalent retrievals as incorrect.

\begin{table*}[tbp]
\centering
\small
\caption{eManual strict, text-equivalent, and deduplicated evaluation}
\label{tab:d1}
\begin{tabularx}{\textwidth}{>{\raggedright\arraybackslash}X>{\raggedright\arraybackslash}X>{\raggedleft\arraybackslash}X>{\raggedleft\arraybackslash}X>{\raggedleft\arraybackslash}X}
\toprule
Method & Evaluation mode & $\mathrm{Hit@10}$ & $\mathrm{MRR@10}$ & $\mathrm{nDCG@10}$ \\
\midrule
BM25 & Strict chunk ID & 0.1154 & 0.0244 & 0.0256 \\
BM25 & Text-equivalent & 0.3846 & 0.3059 & 0.1620 \\
Dense & Strict chunk ID & 0.3231 & 0.0551 & 0.0526 \\
Dense & Text-equivalent & 0.5615 & 0.4716 & 0.2030 \\
Static hybrid & Strict chunk ID & 0.1692 & 0.0366 & 0.0287 \\
Static hybrid & Text-equivalent & 0.5846 & 0.4895 & 0.2263 \\
Dense & Deduplicated corpus & 0.8615 & 0.5736 & 0.3807 \\
\bottomrule
\end{tabularx}
\end{table*}

The text-equivalent and deduplicated metrics are diagnostics, not replacements for the strict evaluation. They show that eManual's low strict score cannot be interpreted as proof that useful local structure is absent.

\section{Encoder Robustness}
\label{sec:12-appendix:encoder-robustness}

The main scale-up uses \nolinkurl{sentence-transformers/all-MiniLM-L6-v2}. A LoTTE 100k robustness check replaces it with \nolinkurl{sentence-transformers/multi-qa-MiniLM-L6-cos-v1}, a QA-tuned MiniLM-family encoder with the same 384-dimensional embedding size and a similar CPU-friendly resource class.

\begin{table*}[tbp]
\centering
\small
\caption{QA-tuned MiniLM-family encoder robustness}
\label{tab:e1}
\begin{tabularx}{\textwidth}{>{\raggedright\arraybackslash}X>{\raggedleft\arraybackslash}X>{\raggedleft\arraybackslash}X>{\raggedleft\arraybackslash}X>{\raggedleft\arraybackslash}X>{\raggedleft\arraybackslash}X>{\raggedleft\arraybackslash}X}
\toprule
Method & $\mathrm{Hit@10}$ & $\mathrm{MRR@10}$ & $\mathrm{nDCG@10}$ & Evid. recall@10 & Avg tokens@10 & Token ratio \\
\midrule
Dense-only & 0.8809 & 0.7220 & 0.6616 & 0.7163 & 1514.51 & 1.0000x \\
Conservative policy & 0.8853 & 0.7118 & 0.6291 & 0.6789 & 1463.71 & 0.9665x \\
\bottomrule
\end{tabularx}
\end{table*}

Under the QA-tuned encoder, the dense baseline becomes stronger and the conservative policy still preserves dense-level query hit while reducing final context tokens by 3.35\%. Ranking metrics and evidence recall are lower than dense, so this is a bounded robustness result rather than a universal retrieval-metric improvement.

\section{Downstream Answer-Quality Check}
\label{sec:12-appendix:downstream-answer-quality-check}

A small downstream answer-quality check compares dense top-10 context with the compressed conservative-policy context on 60 sampled LoTTE 100k queries. The same LLM configuration generates and judges answers in this check.

\begin{table*}[tbp]
\centering
\small
\caption{Downstream answer-quality check}
\label{tab:f1}
\begin{tabularx}{\textwidth}{>{\raggedright\arraybackslash}X>{\raggedleft\arraybackslash}X>{\raggedleft\arraybackslash}X>{\raggedleft\arraybackslash}X>{\raggedleft\arraybackslash}X>{\raggedleft\arraybackslash}X}
\toprule
Method & Answer score & Faithfulness & Answer relevance & Win count & Prompt token ratio \\
\midrule
Dense top-10 & 4.4000 & 4.6500 & 4.6500 & 14 & 1.0000x \\
Conservative policy & 4.2833 & 4.6333 & 4.4500 & 14 & 0.9321x \\
Tie & - & - & - & 32 & - \\
\bottomrule
\end{tabularx}
\end{table*}

The check does not show obvious answer-quality degradation from conservative context compaction. It is not a full human evaluation: the sample is small, one generator/judge model is used, and LLM-as-judge can be biased.

\section{Calibration/Test Context-Budget Validation}
\label{sec:12-appendix:calibration-test-context-budget-validation}

The calibration/test protocol selects the final-context budget on calibration queries and freezes it before evaluation on held-out test queries.

\begin{table*}[tbp]
\centering
\small
\caption{Frozen context-budget validation on LoTTE technology/search}
\label{tab:g1}
\begin{tabularx}{\textwidth}{>{\raggedright\arraybackslash}X>{\raggedright\arraybackslash}X>{\raggedright\arraybackslash}X>{\raggedright\arraybackslash}X>{\raggedleft\arraybackslash}X>{\raggedright\arraybackslash}X>{\raggedleft\arraybackslash}X}
\toprule
Scale & Selected policy & Calibration eligible & Hit delta vs dense & Token saving & Dense trunc. hit delta & Dense trunc. token saving \\
\midrule
100k & \nolinkurl{token_budget_r0.95_m4} & True & +0.00 pp & 6.18\% & -1.44 pp & 13.83\% \\
200k & \nolinkurl{token_budget_r0.85_m4} & True & +1.20 pp & 16.00\% & -2.40 pp & 21.95\% \\
400k & \nolinkurl{token_budget_r0.98_m4} & False / pending follow-up & +2.32 pp & 6.57\% & -0.24 pp & 11.44\% \\
638k & \nolinkurl{token_budget_r0.85_m4} & True & -0.08 pp & 17.53\% & -3.84 pp & 21.90\% \\
\bottomrule
\end{tabularx}
\end{table*}

The calibrated policies should be compared against dense-only adaptive truncation because both reduce final context size. IntentWeight preserves substantially more $\mathrm{Hit@10}$ at a still meaningful token saving level. The 400k row is diagnostic rather than calibration-eligible in the current artifact set and is marked for follow-up calibration.

\section{Cross-Domain Validation}
\label{sec:12-appendix:cross-domain-validation}

LoTTE science/search is used as a second-domain validation, not as a replacement for the main LoTTE technology/search scale-up.

\begin{table*}[tbp]
\centering
\small
\caption{Science/search fixed top-10 ranking validation}
\label{tab:h1}
\begin{tabularx}{\textwidth}{>{\raggedright\arraybackslash}X>{\raggedleft\arraybackslash}X>{\raggedleft\arraybackslash}X>{\raggedleft\arraybackslash}X>{\raggedleft\arraybackslash}X>{\raggedright\arraybackslash}X}
\toprule
Domain/scale & Corpus chunks & Queries & Dense $\mathrm{Hit@10}$ & IntentWeight $\mathrm{Hit@10}$ & Hit delta \\
\midrule
science/search 20k/q200 & 20,490 & 200 & 0.8950 & 0.9267 & +3.17 pp \\
science/search 100k & 101,187 & 596 & 0.8926 & 0.9077 & +1.51 pp \\
\bottomrule
\end{tabularx}
\end{table*}

\begin{table*}[tbp]
\centering
\small
\caption{Science/search frozen context-budget validation}
\label{tab:h2}
\begin{tabularx}{\textwidth}{>{\raggedright\arraybackslash}X>{\raggedright\arraybackslash}X>{\raggedleft\arraybackslash}X>{\raggedright\arraybackslash}X>{\raggedleft\arraybackslash}X>{\raggedright\arraybackslash}X}
\toprule
Domain/scale & Budget policy & Seed & Frozen test hit delta vs dense & Token saving & Strict NI by CI \\
\midrule
20k/q200 & \nolinkurl{token_budget_r0.85_m4} & 13 & +2.86 pp & 13.18\% & True \\
20k/q200 & \nolinkurl{token_budget_r0.85_m4} & 17 & +1.43 pp & 14.31\% & False \\
20k/q200 & \nolinkurl{token_budget_r0.85_m4} & 19 & +0.71 pp & 13.91\% & False \\
100k & \nolinkurl{token_budget_r0.85_m4} & 13 & -1.20 pp & 19.21\% & False \\
100k & \nolinkurl{token_budget_r0.85_m4} & 17 & +0.00 pp & 17.53\% & False \\
100k & \nolinkurl{token_budget_r0.85_m4} & 19 & -0.96 pp & 20.53\% & False \\
\bottomrule
\end{tabularx}
\end{table*}

The fixed top-10 ranking gains transfer, while aggressive final-context budgets require domain and scale calibration.

\section{Feedback-Driven Hard-Case Recovery}
\label{sec:12-appendix:feedback-driven-hard-case-recovery}

Hard-case recovery focuses on affected queries where dense top-10 retrieves at least one GT chunk but the budgeted IntentWeight context misses.

\begin{table*}[tbp]
\centering
\small
\caption{Same-query feedback recovery on affected queries}
\label{tab:i1}
\begin{tabularx}{\textwidth}{>{\raggedright\arraybackslash}X>{\raggedright\arraybackslash}X>{\raggedleft\arraybackslash}X>{\raggedleft\arraybackslash}X>{\raggedleft\arraybackslash}X>{\raggedleft\arraybackslash}X}
\toprule
Domain & Retry method & Affected queries & Recovered & Recovery rate & Avg token saving \\
\midrule
science 100k & arm boost & 34 & 5 & 14.71\% & 17.40\% \\
science 100k & arm boost + conservative budget & 34 & 14 & 41.18\% & 5.76\% \\
science 100k & full-context fallback & 34 & 17 & 50.00\% & -8.07\% \\
technology 100k & arm boost & 42 & 8 & 19.05\% & 13.68\% \\
technology 100k & arm boost + conservative budget & 42 & 9 & 21.43\% & 11.75\% \\
technology 100k & full-context fallback & 42 & 12 & 28.57\% & 0.96\% \\
\bottomrule
\end{tabularx}
\end{table*}

Same-query retry is post-feedback repair evidence. It is not a first-pass generalization result.

\begin{table*}[tbp]
\centering
\small
\caption{Calibration-to-test recovery generalization}
\label{tab:i2}
\begin{tabularx}{\textwidth}{>{\raggedright\arraybackslash}X>{\raggedright\arraybackslash}X>{\raggedright\arraybackslash}X>{\raggedleft\arraybackslash}X}
\toprule
Domain & Recovery policy & Mean hit delta & Avg token saving \\
\midrule
science 100k & conservative budget on learned risky arms after calibration & +0.16 pp & 16.13\% \\
science 100k & full-context fallback on learned risky arms after calibration & +0.48 pp & 13.09\% \\
technology 100k & conservative budget on learned risky arms after calibration & -0.16 pp & 5.88\% \\
technology 100k & full-context fallback on learned risky arms after calibration & +0.16 pp & 4.25\% \\
\bottomrule
\end{tabularx}
\end{table*}

The held-out effect is small and domain-dependent. Feedback should therefore be used as a controlled fallback trigger rather than as unconditional global reranking.

\section{Reproducibility and Reporting Guardrails}
\label{sec:12-appendix:reproducibility-and-reporting-guardrails}

The following rules apply when migrating the draft into a submission template:

\begin{itemize}
\item report query-level $\mathrm{Hit@10}$ as the primary retrieval headline;
\item report $\mathrm{EvidenceRecall@10}$ separately for complete-evidence tasks;
\item use final retrieved context tokens for prompt-context efficiency claims;
\item label source candidate cost and dense invocation rate as retrieval-stage diagnostics;
\item describe multi-epoch prequential adaptation as simulated repeated interaction, not IID held-out generalization;
\item describe feedback as controlled simulation, not collected production feedback;
\item describe geometry diagnostics as support for a piecewise local-structure interpretation, not theorem-level manifold proof;
\item keep dense retrieval visible as a recall floor and fallback route.
\end{itemize}
